\title{: Advancing the Boundaries of Mathematical Reasoning in Open Language Models}

\begin{abstract}
Mathematical reasoning remains a formidable challenge for language models due to its intricate and structured characteristics. In this study, we present DeepSeekMath~7B, which extends the pre-training of DeepSeek-Coder-Base-v1.5 7B by incorporating 120B math-related tokens extracted from Common Crawl, alongside natural language and code datasets. DeepSeekMath~7B attains a remarkable 51.7\% score on the competition-level MATH benchmark without the aid of external tools or ensemble voting methods, nearing the performance levels of Gemini-Ultra and GPT-4. Utilizing self-consistency with 64 samples from DeepSeekMath~7B yields a 60.9\% score on MATH. The enhanced mathematical reasoning capabilities of DeepSeekMath~stem from two main factors: First, leveraging the vast potential of publicly accessible web data through a carefully designed data selection process; second, introducing Group Relative Policy Optimization (GRPO), a variant of Proximal Policy Optimization (PPO), which boosts mathematical reasoning performance while simultaneously improving PPO’s memory efficiency.
\end{abstract}

\section{Introduction}

Large language models (LLMs) have transformed the way mathematical reasoning is approached in artificial intelligence, driving substantial progress on benchmarks such as quantitative reasoning \citep{MATH} and geometry reasoning \citep{trinh2024solving}. Additionally, these models have demonstrated their value in aiding humans to tackle complex mathematical challenges \citep{tao}. Nevertheless, state-of-the-art models like GPT-4 \citep{gpt4} and Gemini-Ultra \citep{gemini} remain inaccessible to the public, and existing open-source models lag considerably behind in performance.

In this work, we present DeepSeekMath, a language model specialized in mathematics that considerably surpasses the capabilities of existing open-source models and approaches the performance of GPT-4 on scholarly benchmarks. To accomplish this, we develop the DeepSeekMath Corpus, a large-scale, high-quality pre-training dataset containing 120B math-related tokens. This corpus is derived from the Common Crawl (CC) using a fastText-based classifier \citep{joulin2016fasttext}. Initially, the classifier is trained with positive samples from OpenWebMath \citep{openwebmath} and a varied set of other web pages as negative samples. Afterward, we apply the classifier to extract more positive samples from CC, which are then further improved through human annotation. The classifier is subsequently retrained with this enhanced dataset to boost its accuracy. Evaluation results demonstrate that the large-scale corpus maintains high quality, with our base model DeepSeekMath-Base 7B attaining 64.2\% on GSM8K \citep{gsm8k} and 36.2\% on the competition-level MATH dataset \citep{MATH}, outperforming Minerva 540B \citep{minerva}. Moreover, since the DeepSeekMath Corpus is multilingual, we observe gains on Chinese mathematical benchmarks \citep{wei2023cmath,agieval}. We consider our efforts in mathematical data curation as a foundation for the research community, recognizing considerable potential for future enhancements.

DeepSeekMath-Base is initialized from DeepSeek-Coder-Base-v1.5 7B \citep{deepseek-coder}, as we found that beginning with a code-trained model is preferable to using a general LLM. Additionally, we observe that math-focused training enhances the model's performance on the MMLU \citep{mmlu} and BBH benchmarks \citep{bbh}, suggesting that it not only improves mathematical skills but also strengthens overall reasoning abilities.

Following pre-training, we conduct mathematical instruction tuning on DeepSeekMath-Base utilizing chain-of-thought \citep{cot}, program-of-thought \citep{pot,pal}, and tool-integrated reasoning \citep{tora} datasets. The resultant model, DeepSeekMath-Instruct 7B, surpasses all other 7B models and achieves performance comparable to 70B-scale open-source instruction-tuned models.

Moreover, we propose Group Relative Policy Optimization (GRPO), a variant of the Proximal Policy Optimization (PPO) \citep{schulman2017proximal} reinforcement learning algorithm. GRPO eliminates the need for a critic model by estimating baselines from group scores, significantly lowering the computational cost of training. Using only a subset of English instruction tuning data, GRPO achieves notable improvements over the strong DeepSeekMath-Instruct model, both on in-domain tasks (GSM8K: 82.9\% $\rightarrow$ 88.2\%, MATH: 46.8\% $\rightarrow$ 51.7\%) and on out-of-domain mathematical tasks (e.g., CMATH: 84.6\% $\rightarrow$ 88.8\%) during the reinforcement learning stage. We also present a unified framework to interpret various methods, including Rejection Sampling Fine-Tuning (RFT) \citep{yuan2023scaling}, Direct Preference Optimization (DPO) \citep{dpo}, PPO, and GRPO. Within this framework, we identify these approaches as either direct or simplified forms of reinforcement learning techniques. Extensive experiments—such as comparisons of online versus offline training, outcome-based versus process-based supervision, and single-turn versus iterative RL—are conducted to thoroughly explore the core components of this framework. Finally, we explain why our RL approach enhances the performance of instruction-tuned models and outline potential avenues for developing more effective RL methods based on this unified perspective.

\subsection{Contributions}
Our work presents scalable mathematics pre-training, along with an in-depth exploration and evaluation of reinforcement learning techniques.

\textbf{Large-Scale Math Pre-Training}

\begin{itemize}[topsep=0pt]
    \item We provide strong evidence that publicly available Common Crawl data holds significant mathematical content. By developing a carefully crafted data selection pipeline, we create the DeepSeekMath Corpus, a high-quality dataset comprising 120B tokens extracted from web pages focused on mathematical material. This dataset is nearly 7 times larger than the math web page corpus used by Minerva \citep{minerva} and 9 times the size of the newly introduced OpenWebMath dataset \citep{openwebmath}.

    \item Our base model, DeepSeekMath-Base 7B, pre-trained on this data, attains performance comparable to Minerva 540B \citep{minerva}, suggesting that model size alone does not determine mathematical reasoning proficiency. Instead, a smaller model trained on high-quality data can also deliver strong results.

    \item We report insights from our math training experiments, showing that pre-training on code before math training enhances the model's capability to solve mathematical tasks both with and without external tool usage. This provides a partial resolution to the persistent question: \textit{does code training enhance reasoning skills?} Our findings indicate that it does, at least within the scope of mathematical reasoning.

    \item Although incorporating arXiv papers in training is a widespread practice, especially in math-focused research, we observe no significant improvement across all mathematical benchmarks considered in this study.
\end{itemize}

\textbf{Investigation and Evaluation of Reinforcement Learning}

\begin{itemize}[topsep=0pt]
    \item We propose Group Relative Policy Optimization (GRPO), a reinforcement learning algorithm that is both efficient and effective. Instead of relying on a critic model, GRPO estimates the baseline using group scores, which substantially lowers the training resource requirements compared to Proximal Policy Optimization (PPO).
    \item Our results show that GRPO considerably improves the performance of our instruction-tuned model, DeepSeekMath-Instruct, using only instruction-tuning data. Additionally, we observe gains in out-of-domain performance throughout the reinforcement learning process.
    \item We offer a unified framework to interpret various approaches, including RFT, DPO, PPO, and GRPO. We also carry out extensive experiments—such as comparing online versus offline training, outcome versus process supervision, and single-turn versus iterative reinforcement learning—to thoroughly analyze the critical components of this framework.
    \item Leveraging this unified framework, we investigate the underlying factors contributing to the success of reinforcement learning and outline several promising avenues for enhancing the effectiveness of reinforcement learning in LLMs.
\end{itemize}

\subsection{Summary of Evaluations and Metrics}

\begin{itemize}[topsep=0pt]
    \item \textbf{English and Chinese Mathematical Reasoning}: We perform extensive evaluations of our models on both English and Chinese benchmarks, spanning math problems from elementary to university levels. English benchmarks include GSM8K \citep{gsm8k}, MATH \citep{MATH}, SAT \citep{llemma}, OCW Courses \citep{minerva}, and MMLU-STEM \citep{mmlu}. Chinese benchmarks comprise MGSM-zh \citep{mgsm}, CMATH \citep{wei2023cmath}, Gaokao-MathCloze \citep{agieval}, and Gaokao-MathQA \citep{agieval}. Our assessments measure the models' capacity to generate fully self-contained textual solutions without tool assistance, as well as their ability to solve problems utilizing Python.

    On English benchmarks, DeepSeekMath-Base competes closely with the closed-source Minerva 540B \citep{minerva}, outperforming all open-source base models (such as Mistral 7B \citep{mistral} and Llemma-34B \citep{llemma}), regardless of whether these models received math-specific pre-training, often by a considerable margin. Particularly, DeepSeekMath-Base excels on Chinese benchmarks, likely because unlike previous works \citep{minerva,llemma}, we incorporate high-quality non-English math pre-training data rather than limiting ourselves to English-only datasets. With mathematical instruction tuning and reinforcement learning, the enhanced DeepSeekMath-Instruct and DeepSeekMath-RL achieve strong results, surpassing 50\% accuracy on the competitive MATH dataset—marking a first in the open-source community.
\end{itemize}

\begin{itemize}
    \item \textbf{Formal Mathematics}: We assess DeepSeekMath-Base by performing the informal-to-formal theorem proving task introduced in \citep{dsp_proof} using miniF2F \citep{minif2f} with Isabelle \citep{isabelle} as the selected proof assistant. DeepSeekMath-Base achieves strong few-shot autoformalization results.
    \item \textbf{Natural Language Understanding, Reasoning, and Code}: To develop a thorough evaluation of the models' overall understanding, reasoning, and programming skills, we test DeepSeekMath-Base on the Massive Multitask Language Understanding (MMLU) benchmark \citep{mmlu}, which includes 57 multiple-choice tasks spanning various disciplines; BIG-Bench Hard (BBH) \citep{bbh}, which contains 23 difficult tasks that mainly demand multi-step reasoning; as well as HumanEval \citep{codex} and MBPP \citep{mbpp}, widely used for assessing code language models. Pre-training on mathematical data enhances both language comprehension and reasoning abilities.
\end{itemize}

\section{Math Pre-Training}

\subsection{Data Collection and Decontamination} In this part, we describe the methodology for building the DeepSeekMath Corpus from Common Crawl. As illustrated in Figure \ref{fig:data_collect}, we outline an iterative procedure that demonstrates a systematic approach to collecting a large-scale mathematical corpus from Common Crawl, beginning with a seed corpus (for example, a small but high-quality dataset related to mathematics). Notably, this method can also be extended to other domains such as programming.

Initially, we select OpenWebMath \citep{openwebmath}, a curated set of high-quality mathematical web documents, as our seed corpus. Using this data, we train a fastText model \citep{joulin2016fasttext} to retrieve additional mathematical web pages resembling OpenWebMath. Specifically, we randomly sample 500,000 examples from the seed corpus as positive training samples and 500,000 Common Crawl web pages as negative samples. Training is conducted using an open-source toolkit\footnote{\url{https://fasttext.cc}}, with vector dimensions set to 256, learning rate at 0.1, a maximum word n-gram length of 3, minimum word occurrences of 3, and 3 epochs. To compress the original Common Crawl, URL-based deduplication and near-deduplication techniques are applied, yielding 40 billion HTML pages. We then use the fastText model to recall mathematical web pages from this deduplicated Common Crawl. To eliminate low-quality mathematical content, pages are ranked based on their fastText scores, and only the highest-ranked ones are retained. The amount of data preserved is evaluated through pre-training on the top 40B, 80B, 120B, and 160B tokens. In the initial iteration, we decide to keep the top 40B tokens.

Following the initial round of data gathering, many mathematical web pages remain uncollected, primarily due to the fastText model being trained on a set of positive examples that lack sufficient variety. To address this, we identify additional mathematical web sources to expand the seed corpus, thereby enhancing the fastText model. Specifically, we first segment the entire Common Crawl into distinct domains, where a domain is defined as web pages sharing the same base URL. For each domain, we compute the proportion of web pages collected during the first iteration. Domains with over 10\% of their pages collected are designated as math-related (e.g., \url{mathoverflow.net}).

Next, we manually annotate URLs containing mathematical content within these identified domains (e.g., \url{mathoverflow.net/questions}). Web pages linked to these URLs that were not collected previously are then incorporated into the seed corpus. This method allows us to acquire additional positive examples, which improves the fastText model’s ability to recall more mathematical data in subsequent iterations. After four rounds of data collection, we accumulate 35.5M mathematical web pages, amounting to 120B tokens. In the fourth iteration, we observe that approximately 98\% of data had already been gathered in the third iteration, prompting us to halt further collection.

To prevent contamination of benchmark datasets, we adopt the filtering approach from \cite{deepseek-coder} to remove web pages containing questions or answers from English mathematical benchmarks such as GSM8K \citep{gsm8k} and MATH \citep{MATH}, as well as Chinese benchmarks like CMATH \citep{wei2023cmath} and AGIEval \citep{agieval}. The filtering rule is as follows: any text segment containing a 10-gram sequence that exactly matches any substring from the evaluation benchmarks is excluded from our math training corpus. For benchmark texts shorter than 10 grams but with a minimum length of 3 grams, we perform exact matching to filter out contaminated pages.

\subsection{Validating the Quality of the DeepSeekMath~Corpus}

We conduct pre-training experiments to compare the DeepSeekMath~Corpus with several recently released math-training datasets:

\begin{itemize}[topsep=0pt]
    \item \textbf{MathPile} \citep{mathpile}: a multi-source corpus (8.9B tokens) compiled from textbooks, Wikipedia, ProofWiki, CommonCrawl, StackExchange, and arXiv, with the majority (over 85\%) derived from arXiv;
    \item \textbf{OpenWebMath} \citep{openwebmath}: CommonCrawl data filtered to retain mathematical content, totaling 13.6B tokens;
    \item \textbf{Proof-Pile-2} \citep{llemma}: a mathematical corpus comprising OpenWebMath, AlgebraicStack (10.3B tokens of mathematical code), and arXiv papers (28.0B tokens). In experiments using Proof-Pile-2, we follow \cite{llemma} in employing an arXiv:Web:Code ratio of 2:4:1.
\end{itemize}

\subsubsection{Training Setup}
\label{sec:quality-policy}
We perform math training on a general pre-trained language model comprising 1.3B parameters, which adopts the same architecture as the DeepSeek LLMs \citep{deepseek-llm}, referred to as DeepSeek-LLM 1.3B. Each mathematical corpus is used to train a separate model for 150B tokens. All experiments are executed using the efficient and lightweight HAI-LLM \citep{haillm} training framework. Following the DeepSeek LLMs training procedure, we utilize the AdamW optimizer \citep{adamW} with $\beta_1=0.9$, $\beta_2=0.95$, and a weight decay of 0.1. The learning rate schedule involves multiple steps: the rate peaks after 2,000 warm-up steps, drops to 31.6\% of its maximum at 80\% of the training, and further decreases to 10.0\% of the peak after 90\% of training completion. The peak learning rate is set to 5.3e-4, and a batch size of 4 million tokens with a context length of 4K is employed.

\subsubsection{Evaluation Outcomes}

\textbf{The DeepSeekMath~Corpus exhibits high quality, multilingual mathematical content, and is the largest available dataset.}

\begin{itemize}[topsep=0pt]
    \item \textbf{High Quality}:
    We assess downstream performance on 8 mathematical benchmarks using few-shot chain-of-thought prompting \cite{cot}. Table \ref{tab:corpora_comparison} clearly shows that models trained on the DeepSeekMath~Corpus outperform others. Figure \ref{fig:corpora_comparisons} illustrates that the model trained on the DeepSeekMath Corpus achieves superior results compared to Proof-Pile-2 after 50B tokens (equivalent to one full epoch of Proof-Pile-2), suggesting that the average quality of the DeepSeekMath Corpus is higher.
    \item \textbf{Multilingual}:
    The DeepSeekMath~Corpus contains data in multiple languages, with English and Chinese being the dominant ones. As demonstrated in Table \ref{tab:corpora_comparison}, training on the DeepSeekMath~Corpus improves mathematical reasoning capabilities in both English and Chinese. On the other hand, existing math corpora, which are mostly English-centric, offer limited gains and may even negatively affect performance in Chinese mathematical reasoning.
    \item \textbf{Large Scale}:
    The DeepSeekMath~Corpus is several times larger than other available math corpora. Figure \ref{fig:corpora_comparisons} shows that DeepSeek-LLM 1.3B, when trained on the DeepSeekMath~Corpus, follows a steeper learning trajectory with more sustained improvements. Conversely, the baseline corpora are smaller and have been reused across multiple training rounds, causing the model performance to plateau rapidly.
\end{itemize}

\subsection{Training and Evaluating DeepSeekMath-Base 7B}

In this section, we present DeepSeekMath-Base 7B, a foundational model exhibiting strong reasoning capabilities, particularly in mathematics. Our model is initialized from DeepSeek-Coder-Base-v1.5 7B \citep{deepseek-coder} and trained on 500B tokens. The data composition is as follows: 56\% derived from the DeepSeekMath Corpus, 4\% from AlgebraicStack, 10\% from arXiv, 20\% from Github code, and the remaining 10\% consists of natural language data from Common Crawl in both English and Chinese. Our primary training configuration aligns with that detailed in Section \ref{sec:quality-policy}, with modifications including a maximum learning rate set at 4.2e-4 and a batch size of 10M tokens.

We perform an extensive evaluation of the mathematical abilities of DeepSeekMath-Base 7B, concentrating on its capacity to generate self-contained mathematical solutions independently, solve problems using external tools, and execute formal theorem proving. Additionally, we offer a broader assessment of the base model’s capabilities, encompassing natural language understanding, reasoning, and programming proficiency.

\paragraph{Mathematical Problem Solving with Step-by-Step Reasoning}

We assess DeepSeekMath-Base’s ability to solve mathematical problems via few-shot chain-of-thought prompting \citep{cot}, across eight benchmarks in both English and Chinese. These benchmarks include quantitative reasoning tasks (such as GSM8K \citep{gsm8k}, MATH \citep{MATH}, and CMATH \citep{wei2023cmath}) as well as multiple-choice problems (including MMLU-STEM \citep{mmlu} and Gaokao-MathQA \citep{agieval}), spanning various areas of mathematics from basic to university-level difficulty.

As indicated in Table \ref{tab:base_math_cot}, DeepSeekMath-Base 7B achieves the highest performance across all eight benchmarks among open-source base models, which include the widely-used general model Mistral 7B \citep{mistral} and the recently introduced Llemma 34B \citep{llemma}, trained on Proof-Pile-2 \citep{llemma} with a focus on math. Remarkably, on the competition-level MATH dataset, DeepSeekMath-Base exceeds other existing open-source base models by more than 10\% absolute margin, and even outperforms Minerva 540B \citep{minerva}, a closed-source base model that is 77 times larger, built upon PaLM \citep{palm} and further trained on mathematical literature.

\paragraph{Mathematical Problem Solving with Tool Use} We assess program-assisted mathematical reasoning on GSM8K and MATH datasets using few-shot program-of-thought prompting \citep{pot,pal}. Models are prompted to address each problem by generating a Python script where libraries such as \textit{math} and \textit{sympy} may be employed for complex calculations. The output of the program execution is taken as the solution. As demonstrated in Table \ref{tab:base_math_pot_proof}, DeepSeekMath-Base 7B surpasses the previous state-of-the-art model, Llemma 34B.

\paragraph{Formal Mathematics}

Automating formal proofs is valuable for verifying the correctness and dependability of mathematical proofs and for boosting efficiency, which has attracted growing interest in recent years. We test DeepSeekMath-Base 7B on the informal-to-formal proving task from \citep{dsp_proof}, which involves generating a formal proof given an informal statement, its formal equivalent, and an informal proof. Our evaluation uses miniF2F \citep{minif2f}, a benchmark for formal Olympiad-level mathematics, where a formal proof in Isabelle is generated for each problem using few-shot prompting. Following \cite{dsp_proof}, we utilize models to produce proof sketches and then employ the pre-existing automated prover Sledgehammer \citep{sledgehammer} to complete the missing components. Table \ref{tab:base_math_pot_proof} shows that DeepSeekMath-Base 7B achieves strong results in proof autoformalization.

\paragraph{Natural Language Understanding, Reasoning, and Code}

We measure model capabilities in natural language understanding on MMLU \citep{mmlu}, reasoning on BBH \citep{bbh}, and coding on HumanEval \citep{codex} and MBPP \citep{mbpp}. According to Table \ref{tab:understanding_reasoning_code}, DeepSeekMath-Base 7B shows considerable improvements over its predecessor, DeepSeek-Coder-Base-v1.5 \citep{deepseek-coder}, in MMLU and BBH, demonstrating the beneficial effects of math training on language understanding and reasoning. Furthermore, by integrating code tokens during continual training, DeepSeekMath-Base 7B successfully preserves the coding performance of DeepSeek-Coder-Base-v1.5 on both coding benchmarks. Overall, DeepSeekMath-Base 7B substantially outperforms the general model Mistral 7B \citep{mistral} on all three reasoning and coding benchmarks.

\section{Supervised Fine-Tuning}

\subsection{SFT Data Preparation}

We develop a mathematical instruction-tuning dataset that encompasses problems in both English and Chinese from various mathematical domains and with different levels of difficulty. Each problem is paired with solutions presented in chain-of-thought (CoT) \citep{cot}, program-of-thought (PoT) \citep{pot,pal}, and tool-enhanced reasoning formats \citep{tora}. The dataset comprises a total of 776K training examples.

\begin{itemize}[topsep=0pt]
    \item \textbf{English mathematical datasets}: We provide tool-integrated solutions for GSM8K and MATH problems, and include a portion of MathInstruct \citep{MathInstruct} alongside the training subset of Lila-OOD \citep{lila}, where problems are addressed using CoT or PoT methods. Our English dataset spans a wide range of mathematical areas such as algebra, probability, number theory, calculus, and geometry.
    \item \textbf{Chinese mathematical datasets}: We compile Chinese K-12 math problems covering 76 distinct sub-topics including linear equations, with solutions annotated in both CoT and tool-enhanced reasoning formats.
\end{itemize}

\subsection{Training and Evaluation of DeepSeekMath-Instruct 7B}

Here, we present DeepSeekMath-Instruct 7B, a model fine-tuned on mathematical instructions starting from DeepSeekMath-Base. Training samples are randomly concatenated until reaching a maximum context window of 4K tokens. The model is trained for 500 steps using a batch size of 256 and a fixed learning rate of 5e-5.

We assess the mathematical capabilities of the models both without and with tool assistance, across four quantitative reasoning benchmarks in English and Chinese. Our model is compared against state-of-the-art models available at the time:

\begin{itemize}[topsep=0pt]
    \item \textbf{Closed-source models} include: (1) the GPT series, with GPT-4 \citep{gpt4} and GPT-4 Code Interpreter~\footnote{\url{https://openai.com/blog/chatgpt-plugins\#\#code-interpreter}} being the most advanced, (2) Gemini Ultra and Pro \citep{gemini}, (3) Inflection-2 \citep{inflection-2}, (4) Grok-1~\footnote{\url{https://x.ai/model-card}}, as well as recently released models from Chinese companies such as (5) Baichuan-3~\footnote{\url{https://www.baichuan-ai.com}}, and (6) the newest GLM-4~\footnote{\url{https://open.bigmodel.cn/dev/api\#glm-4}} from the GLM family \citep{glm}.
\end{itemize}

These models are designed for general use, most of which have undergone multiple alignment processes.
\item \textbf{Open-source models} encompass: general models such as (1) DeepSeek-LLM-Chat 67B \citep{deepseek-llm}, (2) Qwen 72B \citep{qwen}, (3) SeaLLM-v2 7B \citep{seallm}, and (4) ChatGLM3 6B \citep{chatglm3}, along with models enhanced for mathematics including (5) InternLM2-Math 20B~\footnote{\url{https://github.com/InternLM/InternLM-Math}}, which is based on InternLM2 and underwent specialized math training followed by instruction tuning, (6) Math-Shepherd-Mistral 7B which employs PPO training \citep{schulman2017proximal} on Mistral 7B \citep{mistral} using a process-supervised reward model, (7) the WizardMath series \citep{wizardmath} that enhances mathematical reasoning for Mistral 7B and Llama-2 70B \citep{llama2} via evolve-instruct (a variant of instruction tuning with AI-evolved instructions) and PPO training, primarily using training problems from GSM8K and MATH, (8) MetaMath 70B \citep{metamath}, a Llama-2 70B fine-tuned on an expanded version of GSM8K and MATH, (9) ToRA 34B \cite{tora}, a CodeLlama 34B fine-tuned for tool-integrated mathematical reasoning, and (10) MAmmoTH 70B \citep{MathInstruct}, which is Llama-2 70B instruction-tuned on MathInstruct.
\end{itemize}

As detailed in Table \ref{tab:sft_rl_math}, when evaluated without tool usage, DeepSeekMath-Instruct 7B exhibits robust step-by-step reasoning capabilities. Remarkably, on the competition-level MATH dataset, our model outperforms all open-source models and most proprietary models (e.g., Inflection-2 and Gemini Pro) by at least a 9\% absolute margin. This advantage holds even against significantly larger models (e.g., Qwen 72B) or those with specific math-focused reinforcement learning enhancements (e.g., WizardMath-v1.1 7B). Although DeepSeekMath-Instruct competes closely with Chinese proprietary models GLM-4 and Baichuan-3 on MATH, it remains inferior to GPT-4 and Gemini Ultra.

In the evaluation scenario permitting integration of natural language reasoning and program-based tool usage for problem-solving, DeepSeekMath-Instruct 7B achieves nearly 60\% accuracy on MATH, outperforming all current open-source models. On other benchmark datasets, our model remains competitive with DeepSeek-LLM-Chat 67B, the previous state-of-the-art model that is ten times larger.
\section{Reinforcement Learning}

\subsection{Group Relative Policy Optimization} Reinforcement learning (RL) has been demonstrated to effectively enhance the mathematical reasoning capabilities of LLMs following the Supervised Fine-Tuning (SFT) phase \citep{wang2023math,wizardmath}. In this section, we present our efficient and effective RL algorithm, Group Relative Policy Optimization (GRPO).

\subsubsection{From PPO to GRPO} Proximal Policy Optimization (PPO) \citep{schulman2017proximal} is an actor-critic RL method widely employed during the RL fine-tuning stage of LLMs \citep{ouyang2022training}. Specifically, it optimizes LLMs by maximizing the following surrogate objective:

\begin{equation}
\footnotesize
    \mathcal{J}_{PPO}(\theta) = \mathbb{E}_{q \sim P(Q), o \sim \pi_{\theta_{old}}(O|q)} \frac{1}{|o|} \sum_{t=1}^{|o|} \min \left[ \frac{\pi_\theta(o_{t} | q, o_{<t})}{\pi_{\theta_{old}}(o_{t} | q, o_{<t})} A_{t}, \text{clip} \left( \frac{\pi_\theta(o_{t} | q, o_{<t})}{\pi_{\theta_{old}}(o_{t} | q, o_{<t})}, 1 - \varepsilon, 1 + \varepsilon \right)  A_{t} \right] ,
\end{equation}

Here, $\pi_{\theta}$ and $\pi_{\theta_{old}}$ denote the current and previous policy models, while $q$ and $o$ represent questions and outputs sampled from the question dataset and the older policy $\pi_{\theta_{old}}$, respectively. The hyper-parameter $\varepsilon$ is related to clipping and was introduced in PPO to stabilize training. The advantage $A_t$ is calculated via Generalized Advantage Estimation (GAE) \citep{gae}, utilizing the rewards $\{r_{\ge t}\}$ and a learned value function $V_{\psi}$. Therefore, PPO requires training a value function alongside the policy model. To prevent excessive optimization of the reward model, a common method is to incorporate a per-token KL penalty from a reference model into the reward at each token \citep{ouyang2022training}, expressed as

\begin{equation}
    r_{t} = r_\varphi(q, o_{\le t}) - \beta \log\frac{\pi_{\theta}(o_{t}|q, o_{<t})}{\pi_{ref}(o_{t}|q, o_{<t})},
\label{eq:PPO-reward}
\end{equation}

where $r_\varphi$ is the reward model, $\pi_{ref}$ is the reference model—commonly the original SFT model—and $\beta$ represents the weight of the KL penalty.

Since the value function in PPO is typically another model of similar size to the policy model, it imposes a significant demand on memory and computational resources. Moreover, during RL training, the value function acts as a baseline in the advantage calculation to reduce variance. However, in the context of large language models (LLMs), usually only the final token receives a reward score from the reward model, which complicates training a value function that provides accurate token-level estimates. To overcome this, as depicted in Figure \ref{fig:grpo}, we introduce Group Relative Policy Optimization (GRPO), which eliminates the need for an auxiliary value function as required by PPO. Instead, it employs the average reward of multiple outputs sampled in response to the same question as the baseline. Concretely, for each question $q$, GRPO samples a group of outputs $\{o_1, o_2, \cdots, o_G\}$ from the old policy $\pi_{\theta_{old}}$ and updates the policy by maximizing the objective:

\begin{equation}
\footnotesize
\begin{split}
    \mathcal{J}_{GRPO}(\theta) &= \mathbb{E}{[q \sim P(Q), \{o_i\}_{i=1}^G \sim \pi_{\theta_{old}}(O|q)]}  \\
    & \frac{1}{G}\sum_{i=1}^G\frac{1}{|o_i|} \sum_{t=1}^{|o_i|} \left\{ \min \left[ \frac{\pi_\theta(o_{i,t} | q, o_{i,<t})}{\pi_{\theta_{old}}(o_{i,t} | q, o_{i,<t})} \hat{A}_{i,t}, \text{clip} \left( \frac{\pi_\theta(o_{i,t} | q, o_{i,<t})}{\pi_{\theta_{old}}(o_{i,t} | q, o_{i,<t})}, 1 - \varepsilon, 1 + \varepsilon \right)  \hat{A}_{i,t} \right] - \beta \mathbb{D}_{KL}\left[\pi_{\theta} || \pi_{ref}\right]\right\} ,
\end{split}
\label{eq:GRPO-obj}
\end{equation}

where $\varepsilon$ and $\beta$ are hyper-parameters, and $\hat{A}_{i,t}$ is the advantage computed solely from the relative rewards among outputs within each group, to be elaborated in the following subsections. This group-relative calculation of advantages aligns naturally with the comparative nature of reward models, which are generally trained on datasets containing output comparisons for the same question. Furthermore, unlike adding the KL penalty within the reward, GRPO incorporates the KL divergence between the trained policy and the reference policy directly into the loss function, simplifying the computation of $\hat{A}_{i,t}$. In contrast to the KL penalty term in (\ref{eq:PPO-reward}), the KL divergence is estimated using the following unbiased estimator \citep{kl_approx}:

\begin{equation}
\small
    \mathbb{D}_{KL}\left[\pi_{\theta} || \pi_{ref}\right]

\begin{algorithm}[t]
  \small
  \caption{Iterative Group Relative Policy Optimization}
  \textbf{Input} initial policy model $\pi_{\theta_{\text{init}}}$; reward models $r_\varphi$; task prompts $\mathcal{D}$; 
  hyperparameters $\varepsilon$, $\beta$, $\mu$
  \begin{algorithmic}[1]
    \State Initialize policy model $\pi_\theta \leftarrow \pi_{\theta_{\text{init}}}$
    \For{iteration = 1, \dots, I}
       \State Set reference model $\pi_{ref} \leftarrow \pi_{\theta}$
      \For{step = 1, \dots, M}
      \State Draw a batch $\mathcal{D}_b$ from $\mathcal{D}$
      \State Backup the current policy $\pi_{\theta_{old}} \leftarrow \pi_{\theta}$ 
      \State Generate $G$ outputs $\{o_i\}_{i=1}^G \sim \pi_{\theta_{old}} (\cdot \mid q)$ for each query $q \in \mathcal{D}_b$
      \State Evaluate each output $o_i$ using $r_{\varphi}$ to obtain rewards $\{r_i\}_{i=1}^{G}$ 
      \State Calculate $\hat{A}_{i,t}$ for token $t$ of $o_i$ via group relative advantage estimation.
      \For{GRPO iteration = 1, \dots, $\mu$}
        \State Optimize policy $\pi_{\theta}$ by maximizing the GRPO objective (Equation \ref{eq:GC-GRPO})
      \EndFor
    \EndFor \State Update $r_\varphi$ continuously with replay buffer training.
    \EndFor 
  \end{algorithmic}
  \textbf{Output} $\pi_\theta$
  \label{alg:iter-grpo}
\end{algorithm}

\subsubsection{Outcome Supervision RL with GRPO} Specifically, for each question $q$, a set of outputs $\{o_1, o_2, \cdots, o_G\}$ is sampled from the previous policy model $\pi_{\theta_{old}}$. Each output is scored by a reward model, producing $G$ rewards $\mathbf{r}=\{r_1, r_2, \cdots, r_G\}$. These rewards are then normalized by subtracting their mean and dividing by their standard deviation across the group. Outcome supervision assigns the normalized reward at the end of each output $o_i$ and sets the advantage values $\hat{A}_{i, t}$ for all tokens within the output equal to this normalized reward, i.e., $\hat{A}_{i, t} = \widetilde{r}_i = \frac{r_i- {\rm mean}(\mathbf{r})}{{\rm std}(\mathbf{r})}$. The policy is then optimized by maximizing the objective given in equation (\ref{eq:GRPO-obj}).

\subsubsection{Process Supervision RL with GRPO} Outcome supervision only delivers a reward after an entire output, which might not be adequate or efficient for supervising the policy on intricate mathematical problems. Inspired by \cite{wang2023math}, we also investigate process supervision that provides rewards at the conclusion of each reasoning step. Formally, given a question $q$ and $G$ sampled outputs $\{o_1, o_2, \cdots, o_G\}$, a process reward model assigns scores to each step of the outputs, generating rewards: $\mathbf{R} = \{ \{r_1^{index(1)},\cdots,r_1^{index(K_1)}\}, \cdots,  \{r_G^{index(1)},\cdots,r_G^{index(K_G)}\} \}$, where $index(j)$ denotes the token position at the end of the $j$-th step, and $K_i$ represents the total number of steps in output $i$. These stepwise rewards are normalized by subtracting the mean and dividing by the standard deviation over the set $\mathbf{R}$, i.e., $\widetilde{r}_i^{index(j)} = \frac{r_i^{index(j)} - {\rm mean(\mathbf{R})}}{{\rm std(\mathbf{R})}}$. Then, the advantage for each token is computed as the cumulative sum of normalized rewards for all subsequent steps, i.e., $\hat{A}_{i, t} = \sum_{index(j) \ge t} \widetilde{r}_i^{index(j)}$, and the policy is optimized

\subsubsection{Iterative RL with GRPO} As the reinforcement learning training advances, the previously trained reward model may no longer adequately supervise the updated policy model. Consequently, we investigate iterative RL using GRPO. As outlined in Algorithm \ref{alg:iter-grpo}, in iterative GRPO, new training datasets for the reward model are created based on samples drawn from the policy model, and the prior reward model is continuously refined via a replay mechanism that retains 10\% of past data. Subsequently, the policy model is designated as the reference model, and training of the policy model proceeds using the updated reward model.

\subsection{Training and Evaluating DeepSeekMath-RL}

We perform reinforcement learning starting from DeepSeekMath-Instruct 7B. The RL training data consists of chain-of-thought formatted questions derived from GSM8K and MATH within the SFT dataset, totaling approximately 144K questions. Other SFT questions are excluded to specifically assess the effects of RL on benchmarks with limited data during the RL phase. The training set for the reward models is constructed in accordance with \citep{wang2023math}. Our initial reward model is trained based on DeepSeekMath-Base 7B using a learning rate of 2e-5. For GRPO, the policy model’s learning rate is set to 1e-6, and the KL divergence coefficient is 0.04. For each question, we generate $64$ samples. The maximum sequence length is fixed at 1024, and the training batch size is also 1024. The policy model undergoes a single update following each exploration phase. DeepSeekMath-RL 7B is evaluated on benchmarks in the same manner as DeepSeekMath-Instruct 7B. For DeepSeekMath-RL 7B, GSM8K and MATH with chain-of-thought reasoning are considered in-domain tasks, whereas all other benchmarks are treated as out-of-domain tasks.

Table \ref{tab:sft_rl_math} presents the results of open- and closed-source models employing both chain-of-thought and tool-integrated reasoning on English and Chinese evaluation sets. Our observations include: 1) DeepSeekMath-RL 7B achieves accuracy rates of 88.2\% on GSM8K and 51.7\% on MATH using chain-of-thought reasoning. This exceeds the performance of all open-source models ranging from 7B to 70B parameters, as well as most closed-source counterparts. 2) Importantly, DeepSeekMath-RL 7B is trained exclusively on chain-of-thought formatted instruction tuning data from GSM8K and MATH, with DeepSeekMath-Instruct 7B as the starting point. Despite the limited training data scope, it consistently outperforms DeepSeekMath-Instruct 7B across all metrics, demonstrating the benefits of reinforcement learning.

\section{Discussion}
In this section, we discuss the insights gained from our pre-training and reinforcement learning experiments.

\subsection{Insights from Pre-Training}

We begin by sharing our experiences related to pre-training. Unless otherwise noted, training configurations follow those described in Section \ref{sec:quality-policy}. It is important to clarify that, when mentioning the DeepSeekMath Corpus in this section, we refer to an 89-billion-token dataset compiled during the second phase of data collection.

\subsubsection{Code Training Enhances Mathematical Reasoning}

A widely held but unproven belief posits that training on code improves reasoning capabilities. We aim to provide partial evidence supporting this, specifically in the context of mathematics: code training enhances models' capacity for mathematical reasoning both with and without tool assistance.

To investigate the impact of code training on mathematical reasoning, we conducted experiments using two training regimes: two-stage training and one-stage training.

\textbf{Two-Stage Training}

\begin{itemize}[topsep=0pt]
    \item \textbf{Code Training for 400B Tokens $\rightarrow$ Math Training for 150B Tokens}: We first train DeepSeek-LLM 1.3B on 400 billion code tokens, followed by an additional 150 billion math tokens;
    \item \textbf{General Training for 400B Tokens $\rightarrow$ Math Training for 150B Tokens}: As a baseline, we also conduct training using general tokens (sampled from a large-scale general corpus curated by DeepSeek-AI) instead of code tokens in the initial phase, aiming to explore the benefits of code tokens compared to general tokens in enhancing mathematical reasoning.
\end{itemize}

\textbf{One-Stage Training}

\begin{itemize}[topsep=0pt]
    \item \textbf{Math Training for 150B Tokens}: We train DeepSeek-LLM 1.3B exclusively on 150 billion math tokens;
    \item \textbf{Training on a mixture of 400B Code Tokens and 150B Math Tokens}: Since conducting math training after code training tends to degrade coding capabilities, we examine whether a combined training approach mixing code and math tokens in a single stage can still enhance mathematical reasoning while reducing catastrophic forgetting.
\end{itemize}

\paragraph{Results}
Tables \ref{tab:code-math} and \ref{tab:code-math-general-eval} present downstream performance under various training configurations.

Code training enhances program-assisted mathematical reasoning in both two-stage and one-stage training scenarios. As indicated in Table \ref{tab:code-math}, in the two-stage setup, code training alone markedly improves the ability to solve GSM8K and MATH problems via Python, with subsequent math training further boosting performance. Notably, in the one-stage training setting, blending code tokens with math tokens successfully alleviates catastrophic forgetting observed in two-stage training and fosters synergy between coding proficiency (Table \ref{tab:code-math-general-eval}) and program-aided mathematical reasoning (Table \ref{tab:code-math}).

Moreover, code training also benefits mathematical reasoning without tool usage. Within the two-stage framework, initial code training yields moderate improvements and enhances the effectiveness of subsequent math training, ultimately resulting in the highest performance. Conversely, mixing code and math tokens in one-stage training reduces performance on mathematical reasoning without tool use. This may be due to the limited scale of DeepSeek-LLM 1.3B, which might be insufficient to fully assimilate both code and math data simultaneously.

\subsubsection{ArXiv Papers Appear Ineffectual for Enhancing Mathematical Reasoning}

ArXiv papers are frequently incorporated as part of the math pre-training datasets \citep{minerva,gpt-f,llemma,mathpile}. Nonetheless, comprehensive investigations into their influence on mathematical reasoning have been limited. Somewhat surprisingly, our experimental results suggest that arXiv papers do not effectively enhance mathematical reasoning. We test models of varying sizes, including DeepSeek-LLM 1.3B and DeepSeek-Coder-Base-v1.5 7B \citep{deepseek-coder}, utilizing arXiv corpora processed through different pipelines:

\begin{itemize}[topsep=0pt]
    \item \textbf{MathPile} \citep{mathpile}: an 8.9B-token corpus prepared with cleaning and heuristic filtering, with over 85\% composed of scientific arXiv papers;
    \item \textbf{ArXiv-RedPajama} \citep{redpajama}: the complete collection of arXiv LaTeX files stripped of preambles, comments, macros, and bibliographies, amounting to 28.0B tokens.
\end{itemize}

In our tests, we independently train DeepSeek-LLM 1.3B for 150B tokens and DeepSeek-Coder-Base-v1.5 7B for 40B tokens on each arXiv dataset. It appears that arXiv papers do not contribute to enhancements in mathematical reasoning. When trained solely on an arXiv corpus, neither model exhibits significant gains, and some even show declines across several mathematical benchmarks of varying difficulty used in this work. These benchmarks span quantitative reasoning tasks such as GSM8K and MATH (Table \ref{tab:arxiv-cot}), multiple-choice exams like MMLU-STEM (Table \ref{tab:arxiv-cot}), and formal mathematics benchmarks like miniF2F (Table \ref{tab:arxiv_atp}).

However, this finding has caveats and should be considered cautiously. We have not yet explored:

\begin{itemize}[topsep=0pt]
    \item The influence of arXiv tokens on specific math-related tasks excluded from this study, for example, theorem informalization—the process of transforming formal statements or proofs into their informal counterparts;
    \item The effects of arXiv tokens when integrated with other data types;
    \item Whether the advantages of arXiv papers might emerge at larger model scales.
\end{itemize}

Therefore, additional investigation is warranted and is deferred to future research.

\subsection{Insights of Reinforcement Learning}
\subsubsection{Towards a Unified Paradigm} Here, we present a unified framework to analyze various training approaches such as SFT, RFT, DPO, PPO, GRPO, and conduct experiments probing the factors within this unified perspective. In general, the gradient with respect to the parameter $\theta$ for a training method can be expressed as:

\begin{equation}
    \nabla_{\theta}\mathcal{J}_{\textcolor{red}{\mathcal{A}}}(\theta) = \mathbb{E}[\underbrace{(q,o) \sim \textcolor{red}{\mathcal{D}}}_{Data \ Source}]\left( \frac{1}{|o|} \sum_{t=1}^{|o|}  \underbrace{GC_{{\mathcal{A}}}(q, o, t, \textcolor{red}{\pi_{{rf}}})}_{Gradient \ Coefficient}  \nabla_{\theta}\log \pi_{\theta}(o_t | q, o_{<t})\right).
\label{eq:objective}
\end{equation}

There are three principal elements involved: 1) \textit{Data Source $\mathcal{D}$}, which specifies the dataset used for training; 2) \textit{Reward Function $\pi_{{rf}}$}, providing the reward signal during training; 3) \textit{Algorithm $\mathcal{A}$}, which combines the training data and the reward signal to compute the gradient coefficient $GC$, determining the strength of penalty or reinforcement applied to the data. We investigate several notable approaches within this unified framework:

\begin{itemize}
    \item  \textbf{Supervised Fine-tuning (SFT)}: SFT adapts a pretrained model by fine-tuning it on a human-curated SFT dataset.
    \item \textbf{Rejection Sampling Fine-tuning (RFT)}: RFT further fine-tunes the SFT model using filtered outputs generated from SFT model samples given SFT questions, selecting outputs based on answer correctness.
    \item \textbf{Direct Preference Optimization (DPO)}: DPO enhances the SFT model by fine-tuning it on augmented outputs sampled from the SFT model, employing a pairwise DPO loss.
    \item \textbf{Online Rejection Sampling Fine-tuning (Online RFT)}: Unlike RFT, Online RFT starts with the SFT model as the policy and fine-tunes it using augmented outputs sampled dynamically from the current policy model.
    \item \textbf{PPO/GRPO}: PPO/GRPO initializes the policy model with the SFT model and further strengthens it through reinforcement using outputs generated from the live policy model.
\end{itemize}

We provide a summary of the components of these methods in Table \ref{tab:data_policy}. For a more comprehensive derivation process, please see Appendix \ref{app:analysis-rl}.

\paragraph{Observation about Data Source} We categorize the data sources into two types: online sampling and offline sampling. Online sampling means the training data comes from exploration outcomes generated by the current policy model during training, whereas offline sampling means the training data is obtained from the sampling results of the initial SFT model. RFT and DPO use the offline sampling approach, while Online RFT and GRPO utilize the online approach.

As illustrated in Figure \ref{fig:rl-analysis}, we observe that Online RFT substantially outperforms RFT across two benchmarks. Specifically, Online RFT performs similarly to RFT in the initial training phase but achieves a clear advantage in later stages, indicating the benefits of online training. This finding is intuitive since, at the beginning, the actor and SFT models are quite similar, resulting in only slight differences in sampled data. However, as training progresses, data sampled from the actor diverges more significantly, making real-time data sampling more advantageous.

\paragraph{Observation about Gradient Coefficient} The algorithm translates the reward signal into a gradient coefficient to update model parameters. In our experiments, we classify the reward function into two categories: `Rule' and `Model'. The Rule category assesses response quality based on answer correctness, while the Model category involves training a reward model to assign scores to each response. The training data for the reward model is derived from rule-based judgments. Equations \ref{eq:GC-RFT} and \ref{eq:GC-GRPO} reveal a key distinction between GRPO and Online RFT: GRPO uniquely modifies its gradient coefficient according to the reward value provided by the reward model, enabling differentiated reinforcement and penalization depending on the response magnitude. Conversely, Online RFT does not incorporate this feature; it refrains from penalizing incorrect responses and applies uniform reinforcement to all responses with correct answers at the same intensity level.

As illustrated in Figure \ref{fig:rl-analysis}, GRPO outperforms online RFT, underscoring the effectiveness of adjusting positive and negative gradient coefficients. Moreover, GRPO+PS demonstrates better results than GRPO+OS, suggesting the advantages of employing fine-grained, step-aware gradient coefficients. Additionally, we investigate iterative RL by performing two rounds of iteration in our experiments. As depicted in Figure \ref{fig:iter-rl}, iterative RL markedly enhances performance, particularly during the first iteration.

\subsubsection{Why Does RL Work?} In this study, we apply reinforcement learning to a subset of instruction tuning data, achieving substantial performance improvements over the instruction tuning model. To further elucidate why reinforcement learning is effective, we assess the Pass@K and Maj@K accuracy of both the Instruct and RL models across two benchmarks. Figure \ref{fig:maj-pass} reveals that RL improves Maj@K performance but not Pass@K. These results suggest that RL enhances the model’s overall robustness in output distribution; in other words, \textbf{the improvement appears to stem from increasing the likelihood of the correct response appearing within the TopK outputs rather than enhancing core capabilities.} Similarly, \citep{wang2023making} identified a \textbf{misalignment issue} in reasoning tasks with the SFT model, demonstrating that reasoning performance in SFT models can be improved via various preference alignment strategies \citep{yuan2023rrhf,song2023preference,wang2023making}.

\subsubsection{How Can We Achieve More Effective RL?}
\label{sec:effec_rl}
We demonstrate that RL performs well on mathematical reasoning tasks and propose a unified framework to interpret different representative training methods. Within this framework, all approaches are seen as either direct or simplified RL techniques. As outlined in Equation \ref{eq:objective}, three main components are involved: Data Source, Algorithm, and Reward Function. We suggest several potential future directions related to these components.

\paragraph{Data Source} The data source serves as the foundational material for all training methods. Specifically, in RL, the data source refers to unlabeled questions with outputs generated by sampling from the policy model. In this work, we only utilize questions from the instruction tuning phase and employ a basic nucleus sampling method to generate outputs. We believe this might explain why our RL pipeline solely improves Maj@K performance. Future work will explore applying our RL pipeline to out-of-distribution question prompts, alongside \textbf{advanced sampling (decoding) strategies}, such as those based on tree-search techniques \citep{yao2023tree}. Additionally, \textbf{efficient inference methods} \citep{xia-etal-2023-speculative,leviathan2023fast,kwon2023efficient,xia2024unlocking}, which govern the exploration efficiency of policy models, will also play a critical role.

\paragraph{Algorithms} Algorithms utilize the data and reward signals to calculate the gradient coefficient, which is then used to update the model parameters. According to Equation \ref{eq:objective}, current methods generally \textbf{TRUST} the reward function's signal to some degree, using it to either increase or decrease the conditional probability of specific tokens. However, it is not always possible to guarantee the reward signal's reliability, especially in highly complex tasks. For instance, even the PRM800K dataset \citep{lightman2023let}, which has been carefully annotated by expert annotators, still contains around 20\% erroneous annotations\footnote{\url{https://github.com/openai/prm800k/issues/12\#issuecomment-1728491852}}. Given this, we aim to investigate reinforcement learning algorithms that remain robust in the presence of noisy reward signals. We posit that such \textbf{WEAK-TO-STRONG} \citep{burns2023weak} alignment approaches could fundamentally transform learning algorithms.

\paragraph{Reward Function} The reward function serves as the origin of the training signal. In reinforcement learning, the reward function is typically implemented via a neural reward model. We identify three crucial directions for reward models: 1) \textbf{Improving the generalization capacity of the reward model.} The reward model must generalize well to out-of-distribution queries and advanced decoding outputs; otherwise, reinforcement learning might only stabilize the distribution of large language models (LLMs) without truly enhancing their core abilities; 2) \textbf{Representing the uncertainty of the reward model.} Capturing uncertainty could serve as a crucial link between weak reward models and weak-to-strong learning algorithms; 3) \textbf{Developing efficient methods to construct high-quality process reward models} that provide detailed training signals throughout the reasoning process \citep{lightman2023let,wang2023math}.

\section{Conclusion, Limitation, and Future Work} We introduce DeepSeekMath, which surpasses all open-source models on the competition-level MATH benchmark and nearly matches the performance of proprietary models. DeepSeekMath~starts from DeepSeek-Coder-v1.5 7B and undergoes continuous training over 500B tokens, with a large portion of the training data comprising 120B math tokens extracted from Common Crawl. Our comprehensive ablation study indicates that web pages provide substantial opportunities for acquiring high-quality mathematical data, whereas arXiv proves less advantageous than initially anticipated. We propose Group Relative Policy Optimization (GRPO), a variant of Proximal Policy Optimization (PPO), which significantly enhances mathematical reasoning capabilities while consuming less memory. Experimental results demonstrate the effectiveness of GRPO even when DeepSeekMath-Instruct 7B has achieved strong benchmark scores. Additionally, we offer a unified framework to conceptualize a range of methods and highlight several promising avenues for more efficient reinforcement learning.

Despite DeepSeekMath's strong performance on quantitative reasoning benchmarks, its abilities in geometry and theorem proving are relatively weaker compared to closed models. For example, in preliminary testing, the model struggles with problems involving triangles and ellipses, suggesting a potential bias in data selection during pre-training and fine-tuning. Moreover, limited by its model size, DeepSeekMath~performs worse than GPT-4 in few-shot learning scenarios. While GPT-4 shows performance gains with few-shot inputs, DeepSeekMath~exhibits comparable results in both zero-shot and few-shot evaluations. Moving forward, we plan to refine our data selection pipeline to build a higher-quality pre-training corpus. Furthermore, we intend to investigate the promising directions outlined in Section \ref{sec:effec_rl} to enhance the reinforcement learning of LLMs. 

\end{document}